\documentclass[aspectratio=169,10pt]{beamer}

% Shared Beamer setup for Fall 2026 course slides.
\mode<presentation>{
  \usetheme{Madrid}
  \usecolortheme{default}
}

\definecolor{CalPolyGreen}{HTML}{154734}
\definecolor{CalPolyGold}{HTML}{BD8B13}
\definecolor{DeckInk}{HTML}{1F2933}
\definecolor{DeckMuted}{HTML}{5F6B7A}

\setbeamercolor{structure}{fg=CalPolyGreen}
\setbeamercolor{palette primary}{bg=CalPolyGreen,fg=white}
\setbeamercolor{palette secondary}{bg=DeckInk,fg=white}
\setbeamercolor{palette tertiary}{bg=CalPolyGold,fg=DeckInk}
\setbeamercolor{frametitle}{bg=CalPolyGreen,fg=white}
\setbeamercolor{title}{fg=CalPolyGreen}
\setbeamercolor{normal text}{fg=DeckInk,bg=white}
\setbeamercolor{block title}{bg=CalPolyGreen,fg=white}
\setbeamercolor{block body}{bg=black!3,fg=DeckInk}

\setbeamertemplate{navigation symbols}{}
\setbeamertemplate{footline}[frame number]
\setbeamertemplate{itemize item}{\color{CalPolyGold}$\blacktriangleright$}
\setbeamertemplate{itemize subitem}{\color{CalPolyGreen}$\bullet$}

\newcommand{\CourseNumber}{Course}
\newcommand{\CourseTitle}{Course Title}
\newcommand{\LectureNumber}{0}
\newcommand{\LectureTitle}{Lecture Title}
\newcommand{\LectureDate}{Date}
\newcommand{\CourseUnit}{Unit}

\newcommand{\coursenumber}[1]{\renewcommand{\CourseNumber}{#1}}
\newcommand{\coursetitle}[1]{\renewcommand{\CourseTitle}{#1}}
\newcommand{\lecturenumber}[1]{\renewcommand{\LectureNumber}{#1}}
\newcommand{\lecturetitle}[1]{\renewcommand{\LectureTitle}{#1}}
\newcommand{\lecturedate}[1]{\renewcommand{\LectureDate}{#1}}
\newcommand{\courseunit}[1]{\renewcommand{\CourseUnit}{#1}}

\author{Kirk Duran}
\institute{California Polytechnic State University, San Luis Obispo}
\date{\LectureDate}

\newcommand{\makedecktitle}{
  \begin{frame}[plain]
    \vspace{0.25in}
    {\usebeamerfont{title}\color{CalPolyGreen}\LectureTitle\par}
    \vspace{0.18in}
    {\large \CourseNumber{}: \CourseTitle\par}
    \vspace{0.08in}
    {\large Lecture \LectureNumber{} -- \LectureDate\par}
    \vfill
    {\small Kirk Duran \textbar{} Cal Poly, San Luis Obispo\par}
  \end{frame}
}

\newcommand{\placeholdernote}{
  \begin{block}{Placeholder}
    Replace these bullets with the final lecture material, examples, and in-class activities.
  \end{block}
}

\AtBeginSection[]{
  \begin{frame}{Roadmap}
    \tableofcontents[currentsection]
  \end{frame}
}


\usepackage{tikz}

\graphicspath{{../assets/lecture-07/}}

% If an image file is not yet available, skip the visual slot.
\newcommand{\lectureimage}[2][1.75in]{%
  \IfFileExists{../assets/lecture-07/#2}{%
    \includegraphics[width=\linewidth,height=#1,keepaspectratio]{#2}%
  }{%
  }%
}

\setbeamersize{text margin left=0.42in,text margin right=0.42in}

\coursenumber{CS 4880}
\coursetitle{Artificial Intelligence}
\lecturenumber{7}
\lecturetitle{Local Search}
\lecturedate{September 9, 2026}
\courseunit{Search and Problem Solving}

\begin{document}

\makedecktitle

\begin{frame}[shrink=20]{Today}
  \begin{itemize}
    \item Why path-based tree search can be the wrong abstraction for large optimization problems.
    \item How local search represents a problem using a complete state, an objective function, and a neighborhood.
    \item How hill climbing improves a single current state and why it can become trapped.
    \item How random restarts and simulated annealing introduce exploration.
    \item How local beam search and genetic algorithms search with populations of candidate states.
    \item Why local search remains useful in classical AI, combinatorial optimization, robotics, and games.
  \end{itemize}

  \begin{keyidea}
    Local search gives up explicit path construction in order to search enormous state spaces with very little memory.
  \end{keyidea}
\end{frame}

\sectionframe{Part 1: When Paths Stop Mattering}

\begin{frame}[shrink=20]{From Path Finding to State Optimization}
  \begin{columns}[T,onlytextwidth]
    \begin{column}{0.55\textwidth}
      Systematic search asks which path connects an initial state to a goal.

      \begin{itemize}
        \item DFS explores deeply before backtracking.
        \item BFS explores level by level.
        \item A\* prioritizes states using path cost and heuristic guidance.
        \item Each method maintains enough information to reconstruct a path.
      \end{itemize}

      \begin{block}{A different question}
        What if we care primarily about the \emph{quality of the final state}, not the sequence of actions used to reach it?
      \end{block}
    \end{column}

    \begin{column}{0.39\textwidth}
      \lectureimage[2.55in]{slide4.png}
    \end{column}
  \end{columns}

  \begin{keyidea}
    Optimization problems often need a good final configuration rather than an explicitly stored route through the state space.
  \end{keyidea}
\end{frame}

\begin{frame}[shrink=20]{When Tree Search Fails}
  \begin{columns}[T,onlytextwidth]
    \begin{column}{0.54\textwidth}
      Tree search can become inefficient or infeasible when:

      \begin{itemize}
        \item the state space is vast or effectively infinite;
        \item only the final state matters;
        \item the problem is naturally phrased as optimization rather than navigation;
        \item storing visited states, predecessors, and complete paths consumes too much memory.
      \end{itemize}

      \begin{block}{Rubik's Cube illustration}
        If each position has roughly 18 legal moves, a naive tree has about 18 nodes at depth 1, 324 at depth 2, and 5,832 at depth 3.
      \end{block}
    \end{column}

    \begin{column}{0.40\textwidth}
      \lectureimage[2.65in]{slide5.png}
    \end{column}
  \end{columns}

  \begin{keyidea}
    Exponential growth makes explicit search trees unattractive when the path itself provides little value.
  \end{keyidea}
\end{frame}

\begin{frame}[shrink=20]{Local Search}
  \begin{cpdefinition}
    \textbf{Local search} operates from one or more current states and improves them using neighboring states, without explicitly storing complete paths through the search space.
  \end{cpdefinition}

  \begin{columns}[T,onlytextwidth]
    \begin{column}{0.52\textwidth}
      \begin{itemize}
        \item Evaluate the current state.
        \item Generate or sample neighboring states.
        \item Move according to an improvement or exploration rule.
        \item Repeat until a solution or stopping condition is reached.
      \end{itemize}

      \begin{block}{Why use it?}
        Local search uses very little memory and can find good solutions even when systematic search is impractical.
      \end{block}
    \end{column}

    \begin{column}{0.42\textwidth}
      \lectureimage[2.45in]{slide6.png}
    \end{column}
  \end{columns}

  \begin{keyidea}
    Local search is not systematic: it may never visit large portions of the state space, including regions containing solutions.
  \end{keyidea}
\end{frame}

\begin{frame}[shrink=20]{A New Problem: Eight Queens}
  \begin{columns}[T,onlytextwidth]
    \begin{column}{0.50\textwidth}
      Place eight queens on an $8\times 8$ chessboard so that no two queens attack one another.

      \begin{block}{Complete-state formulation}
        Keep exactly one queen in each column. A state specifies the row occupied by the queen in every column.
      \end{block}

      \begin{itemize}
        \item Every candidate already contains all eight queens.
        \item Most candidates are invalid because queens attack one another.
        \item Search improves the arrangement until all conflicts disappear.
      \end{itemize}
    \end{column}

    \begin{column}{0.44\textwidth}
      \lectureimage[2.65in]{slide7.png}
    \end{column}
  \end{columns}

  \begin{keyidea}
    Local search often starts with a complete but imperfect solution and iteratively improves it.
  \end{keyidea}
\end{frame}

\sectionframe{Part 2: Objective Functions and State-Space Landscapes}

\begin{frame}[shrink=20]{Objective Function}
  \begin{cpdefinition}
    An \textbf{objective function} assigns a real-valued quality score to a complete candidate state:
    \[
      f(s)\in\mathbb{R}.
    \]
  \end{cpdefinition}

  \begin{columns}[T,onlytextwidth]
    \begin{column}{0.52\textwidth}
      \begin{itemize}
        \item A local-search problem may seek to \textbf{minimize} or \textbf{maximize} this value.
        \item The score provides graded feedback even when the current state is not yet a solution.
        \item The neighborhood function determines which states can be reached by one local move.
      \end{itemize}

      \begin{block}{Optimization model}
        State space $+$ objective function $+$ neighborhood function.
      \end{block}
    \end{column}

    \begin{column}{0.42\textwidth}
      \lectureimage[2.45in]{slide9.png}
    \end{column}
  \end{columns}
\end{frame}

\begin{frame}[shrink=20]{Eight Queens as Optimization}
  \begin{columns}[T,onlytextwidth]
    \begin{column}{0.50\textwidth}
      For the complete-state formulation:

      \begin{itemize}
        \item \textbf{State}: eight queens, one queen per column.
        \item \textbf{Objective}: number of attacking pairs of queens.
        \item \textbf{Goal}: minimize the objective to $0$.
        \item \textbf{Neighbor}: move one queen to a different row in its column.
      \end{itemize}

      Each of 8 queens can move to 7 other rows, giving
      \[
        8\cdot 7=56
      \]
      immediate neighboring states.
    \end{column}

    \begin{column}{0.44\textwidth}
      \lectureimage[2.60in]{slide10.png}
    \end{column}
  \end{columns}

  \begin{keyidea}
    The neighborhood defines the algorithm's local view of the optimization landscape.
  \end{keyidea}
\end{frame}

\begin{frame}[shrink=20]{State-Space Landscape}
  \begin{columns}[T,onlytextwidth]
    \begin{column}{0.53\textwidth}
      A \textbf{state-space landscape} visualizes candidate states according to their objective values.

      \begin{itemize}
        \item Each point represents a complete state.
        \item Height represents objective value.
        \item Neighboring points correspond to locally reachable states.
        \item Local search moves across this landscape by choosing neighboring states.
      \end{itemize}

      \begin{block}{Direction depends on the problem}
        For maximization, ``uphill'' is better. For minimization, the same picture can be interpreted by reversing the objective or descending toward lower cost.
      \end{block}
    \end{column}

    \begin{column}{0.41\textwidth}
      \lectureimage[2.65in]{slide11.png}
    \end{column}
  \end{columns}

  \begin{keyidea}
    Local-search performance depends strongly on the topography induced by the objective and neighborhood functions.
  \end{keyidea}
\end{frame}

\sectionframe{Part 3: Hill Climbing --- Repeated Local Improvement}

\begin{frame}[shrink=20]{Hill Climbing}
  \begin{cpdefinition}
    \textbf{Hill climbing} repeatedly moves from the current state to its best improving neighbor.
  \end{cpdefinition}

  \begin{columns}[T,onlytextwidth]
    \begin{column}{0.53\textwidth}
      \begin{itemize}
        \item Keep only a single current state.
        \item Evaluate the immediate neighborhood.
        \item Move in the direction of greatest improvement.
        \item Stop when no neighboring state improves the objective.
      \end{itemize}

      \begin{block}{Memory}
        Hill climbing does not need to store the path, previously visited states, or a large frontier.
      \end{block}
    \end{column}

    \begin{column}{0.41\textwidth}
      \lectureimage[2.55in]{slide13.png}
    \end{column}
  \end{columns}

  \begin{keyidea}
    Hill climbing is a greedy local optimizer: it uses immediate improvement without planning several moves ahead.
  \end{keyidea}
\end{frame}

\begin{frame}[shrink=20]{Hill-Climbing Pseudocode}
  \begin{block}{Algorithm}
  \ttfamily\small
  current $\leftarrow$ initial\_state\\[0.03in]
  while true:\\
  \quad neighbors $\leftarrow$ NEIGHBORS(current)\\
  \quad next $\leftarrow$ best state in neighbors\\[0.03in]
  \quad if next does not improve current:\\
  \qquad return current\\[0.03in]
  \quad current $\leftarrow$ next
  \end{block}

  \begin{itemize}
    \item For maximization, ``best'' means highest objective value.
    \item For minimization, ``best'' means lowest objective value.
    \item The returned state may be a global optimum, a local optimum, or a plateau state.
  \end{itemize}

  \begin{keyidea}
    The algorithm's simplicity is also its weakness: a one-step local decision cannot see a better region separated by worse intermediate states.
  \end{keyidea}
\end{frame}

\begin{frame}[shrink=20]{Hill Climbing on Eight Queens}
  \begin{columns}[T,onlytextwidth]
    \begin{column}{0.52\textwidth}
      Use the number of attacking pairs as the cost $h(s)$.

      \begin{itemize}
        \item Start with eight queens, one per column.
        \item Generate the 56 states obtained by moving one queen within its column.
        \item Choose the neighbor with the fewest attacking pairs.
        \item Repeat until $h(s)=0$ or no improvement exists.
      \end{itemize}

      \begin{block}{Goal state}
        \[
          h(s)=0
        \]
        means that no pair of queens attacks one another.
      \end{block}
    \end{column}

    \begin{column}{0.42\textwidth}
      \lectureimage[2.60in]{slide15.png}
    \end{column}
  \end{columns}
\end{frame}

\begin{frame}[shrink=20]{Why Hill Climbing Gets Stuck}
  \begin{columns}[T,onlytextwidth]
    \begin{column}{0.53\textwidth}
      The state-space landscape can contain structures that defeat purely greedy improvement.

      \begin{itemize}
        \item \textbf{Local maximum}: better than every neighbor but worse than a distant state.
        \item \textbf{Ridge}: improvement requires coordinated movement not captured by one simple local step.
        \item \textbf{Plateau}: a flat region in which many neighboring states have equal value.
        \item \textbf{Shoulder}: a flat region from which an improving move eventually becomes available.
      \end{itemize}
    \end{column}

    \begin{column}{0.41\textwidth}
      \lectureimage[2.70in]{slide16.png}
    \end{column}
  \end{columns}

  \begin{keyidea}
    A locally optimal state need not be globally optimal.
  \end{keyidea}
\end{frame}

\begin{frame}[shrink=20]{Sideways Moves}
  \begin{columns}[T,onlytextwidth]
    \begin{column}{0.54\textwidth}
      A \textbf{sideways move} transitions to a neighboring state with the same objective value.

      \begin{itemize}
        \item It neither improves nor worsens the current solution.
        \item It can move across a plateau until an improving exit is found.
        \item It can also wander indefinitely through a flat region.
      \end{itemize}

      \begin{block}{Practical safeguard}
        Limit the number of consecutive sideways moves; the source example suggests a cap such as 100.
      \end{block}
    \end{column}

    \begin{column}{0.40\textwidth}
      \lectureimage[2.55in]{slide17.png}
    \end{column}
  \end{columns}

  \begin{keyidea}
    Allowing neutral moves expands local mobility, but it requires a stopping rule to avoid infinite wandering.
  \end{keyidea}
\end{frame}

\begin{frame}[shrink=20]{Stochastic Hill Climbing}
  \begin{columns}[T,onlytextwidth]
    \begin{column}{0.53\textwidth}
      Instead of always selecting the single steepest improvement:

      \begin{itemize}
        \item identify neighbors that improve the objective;
        \item choose randomly among the available improving moves;
        \item repeat from the newly selected state.
      \end{itemize}

      \begin{block}{Tradeoff}
        Randomness can explore different uphill trajectories and reduce deterministic trapping, but convergence may be slower.
      \end{block}
    \end{column}

    \begin{column}{0.41\textwidth}
      \lectureimage[2.55in]{slide18.png}
    \end{column}
  \end{columns}

  \begin{keyidea}
    Random choice among improving moves creates diversity without yet accepting worse states.
  \end{keyidea}
\end{frame}

\begin{frame}[shrink=20]{Best-First Hill Climbing}
  \begin{columns}[T,onlytextwidth]
    \begin{column}{0.54\textwidth}
      The source lecture describes a variant that examines successors incrementally rather than evaluating the full neighborhood.

      \begin{itemize}
        \item Generate candidate successors one at a time.
        \item Stop when the first improving neighbor is found.
        \item Move immediately to that state.
        \item Avoid scoring every successor when neighborhoods are large.
      \end{itemize}

      \begin{block}{When it helps}
        This approach can reduce evaluation cost when there are many possible successors and improvements are common.
      \end{block}
    \end{column}

    \begin{column}{0.40\textwidth}
      \lectureimage[2.50in]{slide19.png}
    \end{column}
  \end{columns}
\end{frame}

\begin{frame}[shrink=20]{Random-Restart Hill Climbing}
  \begin{columns}[T,onlytextwidth]
    \begin{column}{0.52\textwidth}
      Random restart runs hill climbing repeatedly from independently sampled initial states.

      \begin{itemize}
        \item Each run explores a different region of the landscape.
        \item A failed run does not force the next run to begin from the same local optimum.
        \item Continue until a satisfactory solution is found or a restart budget is exhausted.
      \end{itemize}

      \begin{block}{Effective when}
        The state space contains many local optima and finding a sufficiently good solution matters more than following one deterministic trajectory.
      \end{block}
    \end{column}

    \begin{column}{0.42\textwidth}
      \lectureimage[2.65in]{slide20.png}
    \end{column}
  \end{columns}

  \begin{keyidea}
    Restarts trade additional time for robustness across difficult landscapes.
  \end{keyidea}
\end{frame}

\sectionframe{Part 4: Simulated Annealing --- Sometimes Move the Wrong Way}

\begin{frame}[shrink=20]{Simulated Annealing}
  \begin{cpdefinition}
    \textbf{Simulated annealing} combines local improvement with controlled random exploration by sometimes accepting a worse neighboring state.
  \end{cpdefinition}

  \begin{columns}[T,onlytextwidth]
    \begin{column}{0.52\textwidth}
      At each step:

      \begin{itemize}
        \item select a random neighboring state;
        \item accept it if it improves the objective;
        \item if it is worse, accept it with a probability controlled by temperature;
        \item gradually reduce that temperature over time.
      \end{itemize}

      \begin{block}{Physical analogy}
        Metallurgical annealing uses high heat for random motion, followed by gradual cooling toward a stable crystal structure.
      \end{block}
    \end{column}

    \begin{column}{0.42\textwidth}
      \lectureimage[2.65in]{slide22.png}
    \end{column}
  \end{columns}
\end{frame}

\begin{frame}[shrink=20]{Simulated-Annealing Pseudocode}
  \begin{block}{Algorithmic structure}
  \ttfamily\small
  current $\leftarrow$ initial\_state\\[0.03in]
  for step $=1,2,\ldots$:\\
  \quad $T \leftarrow$ TEMPERATURE(step)\\
  \quad if $T$ is sufficiently low: return current\\[0.03in]
  \quad next $\leftarrow$ random neighbor of current\\
  \quad if next improves current:\\
  \qquad current $\leftarrow$ next\\
  \quad else if ACCEPT\_WORSE(next,current,$T$):\\
  \qquad current $\leftarrow$ next
  \end{block}

  \begin{keyidea}
    Early randomness can escape local traps; later low temperature makes the search increasingly conservative.
  \end{keyidea}
\end{frame}

\begin{frame}[shrink=20]{Temperature Controls Exploration}
  \begin{columns}[T,onlytextwidth]
    \begin{column}{0.47\textwidth}
      \begin{block}{High temperature}
        \begin{itemize}
          \item high randomness;
          \item worse moves are accepted more often;
          \item broad exploration of the state space.
        \end{itemize}
      \end{block}

      \begin{block}{Low temperature}
        \begin{itemize}
          \item low randomness;
          \item worse moves are rarely accepted;
          \item search concentrates near promising states.
        \end{itemize}
      \end{block}
    \end{column}

    \begin{column}{0.47\textwidth}
      \begin{block}{Cooling schedule}
        \begin{itemize}
          \item fast cooling: quicker convergence, greater risk of getting stuck;
          \item slow cooling: more exploration and computation, greater opportunity to reach a better optimum.
        \end{itemize}
      \end{block}

      \lectureimage[1.40in]{slide24.jpeg}
    \end{column}
  \end{columns}

  \begin{keyidea}
    The cooling schedule determines how quickly simulated annealing shifts from exploration to exploitation.
  \end{keyidea}
\end{frame}

\sectionframe{Part 5: Population-Based Local Search}

\begin{frame}[shrink=20]{Local Beam Search}
  \begin{cpdefinition}
    \textbf{Local beam search} maintains $k$ candidate states in parallel instead of a single current state.
  \end{cpdefinition}

  \begin{columns}[T,onlytextwidth]
    \begin{column}{0.53\textwidth}
      Each iteration:

      \begin{itemize}
        \item generate successors from all $k$ current states;
        \item evaluate the combined successor pool;
        \item retain the best $k$ states for the next iteration.
      \end{itemize}

      \begin{block}{Beam width}
        $k=1$ behaves like single-state hill climbing. Larger $k$ increases breadth and exploration, but also computation and memory.
      \end{block}
    \end{column}

    \begin{column}{0.41\textwidth}
      \lectureimage[2.60in]{slide26.png}
    \end{column}
  \end{columns}

  \begin{keyidea}
    A population lets several regions of the state space compete at the same time.
  \end{keyidea}
\end{frame}

\begin{frame}[shrink=20]{Beam Search Can Lose Diversity}
  \begin{columns}[T,onlytextwidth]
    \begin{column}{0.48\textwidth}
      \begin{block}{Premature convergence}
        If the best candidates become too similar, the beam may collapse into one region of the state space.
      \end{block}

      \begin{block}{Mitigations from the source}
        \begin{itemize}
          \item begin from diverse initial states;
          \item use a sufficiently wide beam;
          \item preserve enough variation for competing regions to survive.
        \end{itemize}
      \end{block}
    \end{column}

    \begin{column}{0.46\textwidth}
    \end{column}
  \end{columns}
\end{frame}

\begin{frame}[shrink=20]{The ``Hive Mind'' Intuition}
  \begin{columns}[T,onlytextwidth]
    \begin{column}{0.48\textwidth}
      The source motivates population search using four ideas:

      \begin{enumerate}
        \item \textbf{Crowd}: generate many candidate solutions.
        \item \textbf{Filter}: keep the candidates with the strongest fitness.
        \item \textbf{Synthesis}: combine useful pieces from different candidates.
        \item \textbf{Spark}: introduce a small random change to avoid a rut.
      \end{enumerate}
    \end{column}

    \begin{column}{0.46\textwidth}
    \end{column}
  \end{columns}

  \begin{keyidea}
    Population methods replace one trajectory with repeated selection, recombination, and diversity.
  \end{keyidea}
\end{frame}

\begin{frame}[shrink=20]{Genetic Algorithms}
  \begin{cpdefinition}
    A \textbf{genetic algorithm} maintains a population of candidate solutions and repeatedly applies fitness-based selection, crossover, and mutation.
  \end{cpdefinition}

  \begin{columns}[T,onlytextwidth]
    \begin{column}{0.30\textwidth}
      \begin{block}{Selection}
        Candidates with higher fitness are more likely to reproduce.
      \end{block}
    \end{column}

    \begin{column}{0.32\textwidth}
      \begin{block}{Crossover}
        Two parent candidates combine portions of their representations to produce offspring.
      \end{block}
    \end{column}

    \begin{column}{0.32\textwidth}
      \begin{block}{Mutation}
        Small random changes introduce new variation into the population.
      \end{block}
    \end{column}
  \end{columns}

  \vspace{0.10in}
  \lectureimage[1.45in]{slide29.png}

  \begin{keyidea}
    Genetic algorithms perform global exploration by evolving a population rather than improving only one current state.
  \end{keyidea}
\end{frame}

\begin{frame}[shrink=20]{Genetic-Algorithm Pseudocode}
  \begin{columns}[T,onlytextwidth]
    \begin{column}{0.60\textwidth}
      \begin{block}{Algorithmic structure}
      \ttfamily\scriptsize
      population $\leftarrow$ RANDOM\_POPULATION()\\[0.03in]
      while stopping condition not met:\\
      \quad evaluate FITNESS for every individual\\
      \quad new\_population $\leftarrow [\ ]$\\[0.03in]
      \quad while new\_population is not full:\\
      \qquad parent1, parent2 $\leftarrow$ SELECT(population)\\
      \qquad child $\leftarrow$ CROSSOVER(parent1,parent2)\\
      \qquad child $\leftarrow$ MUTATE(child)\\
      \qquad add child to new\_population\\[0.03in]
      \quad population $\leftarrow$ new\_population\\[0.03in]
      return best individual
      \end{block}
    \end{column}

    \begin{column}{0.34\textwidth}
      \lectureimage[2.65in]{slide30.jpeg}
    \end{column}
  \end{columns}

  \begin{keyidea}
    Selection increases pressure toward strong solutions, while crossover and mutation generate new combinations for later generations.
  \end{keyidea}
\end{frame}

\begin{frame}[shrink=20]{Genetic Algorithms on Eight Queens}
  \begin{columns}[T,onlytextwidth]
    \begin{column}{0.54\textwidth}
      The source encodes an individual as eight row indices:
      \[
        [r_0,r_1,\ldots,r_7],
      \]
      where index $i$ is a column and $r_i$ is the queen's row.

      \begin{itemize}
        \item \textbf{Initial population}: 100 random individuals.
        \item \textbf{Fitness}: number of non-attacking queen pairs.
        \item \textbf{Maximum fitness}: $\binom{8}{2}=28$.
        \item \textbf{Selection}: tournament among 3 randomly selected individuals.
        \item \textbf{Crossover}: combine two parents at a random point.
        \item \textbf{Mutation}: with small probability (source example: $0.05$), change one queen's row.
      \end{itemize}
    \end{column}

    \begin{column}{0.40\textwidth}
      \lectureimage[2.75in]{slide31.png}
    \end{column}
  \end{columns}
\end{frame}

\begin{frame}[shrink=20]{What Genetic Algorithms Trade}
  \begin{columns}[T,onlytextwidth]
    \begin{column}{0.47\textwidth}
      \begin{block}{Strengths}
        \begin{itemize}
          \item highly parallel population search;
          \item effective global exploration;
          \item useful when direct gradient information is unavailable or misleading;
          \item recombination can preserve useful traits from multiple candidates.
        \end{itemize}
      \end{block}
    \end{column}

    \begin{column}{0.47\textwidth}
      \begin{block}{Limitations}
        \begin{itemize}
          \item no general guarantee of convergence;
          \item performance depends on population size;
          \item crossover and mutation rates require tuning;
          \item populations can still lose diversity.
        \end{itemize}
      \end{block}
    \end{column}
  \end{columns}

  \begin{keyidea}
    Population diversity is valuable only if selection pressure does not eliminate it too quickly.
  \end{keyidea}
\end{frame}

\sectionframe{Part 6: Comparing Local-Search Strategies}

\begin{frame}[shrink=20]{Local Search at a Glance}
  \begin{center}
  \renewcommand{\arraystretch}{1.30}
  \begin{tabular}{l|c|l}
    \textbf{Algorithm} & \textbf{Current states} & \textbf{Main exploration mechanism} \\\hline
    Hill climbing & 1 & best local improvement \\
    Stochastic hill climbing & 1 & random improving move \\
    Random restart & 1 per run & new random initial states \\
    Simulated annealing & 1 & occasionally accept worse moves \\
    Local beam search & $k$ & retain best states from a population \\
    Genetic algorithm & population & selection, crossover, and mutation
  \end{tabular}
  \end{center}

  \begin{keyidea}
    These algorithms differ primarily in how much diversity they preserve and what kinds of non-greedy exploration they permit.
  \end{keyidea}
\end{frame}

\begin{frame}[shrink=20]{Exploration vs. Exploitation}
  \begin{columns}[T,onlytextwidth]
    \begin{column}{0.47\textwidth}
      \begin{block}{Exploitation}
        Prefer states that already appear better.
        \begin{itemize}
          \item hill climbing;
          \item low-temperature annealing;
          \item strong fitness selection.
        \end{itemize}
      \end{block}
    \end{column}

    \begin{column}{0.47\textwidth}
      \begin{block}{Exploration}
        Deliberately preserve alternatives or tolerate temporary regressions.
        \begin{itemize}
          \item stochastic choices;
          \item random restarts;
          \item high-temperature annealing;
          \item wider beams and mutation.
        \end{itemize}
      \end{block}
    \end{column}
  \end{columns}

  \begin{keyidea}
    Local-search design is largely the problem of balancing rapid improvement against the risk of premature convergence.
  \end{keyidea}
\end{frame}

\begin{frame}[shrink=20]{Is This AI?}
  \begin{columns}[T,onlytextwidth]
    \begin{column}{0.54\textwidth}
      Local search belongs to the classical AI toolkit.

      \begin{itemize}
        \item Like BFS, DFS, and A\*, it searches a space of possible solutions.
        \item Unlike path-finding algorithms, it emphasizes state quality rather than reconstructing a complete action sequence.
        \item Its mechanisms may feel systematic compared with modern generative AI, but the underlying search problems remain central to AI.
        \item The source highlights continued relevance in combinatorial optimization, robotics, and game AI.
      \end{itemize}
    \end{column}

    \begin{column}{0.40\textwidth}
    \end{column}
  \end{columns}

  \begin{keyidea}
    AI includes both learned models and explicit search procedures for finding useful solutions in difficult state spaces.
  \end{keyidea}
\end{frame}

\begin{frame}[shrink=20]{Takeaways}
  \begin{itemize}
    \item Local search focuses on the quality of complete states rather than storing explicit paths through a search tree.
    \item An optimization formulation requires a state space, an objective function, and a neighborhood function.
    \item Hill climbing is memory efficient but can become trapped by local optima, ridges, and plateaus.
    \item Sideways moves, stochastic choices, and random restarts increase robustness while retaining a single-state search model.
    \item Simulated annealing accepts controlled temporary regressions so that the search can escape local traps.
    \item Beam search and genetic algorithms maintain populations, using diversity to explore multiple regions at once.
    \item Next lecture: \textbf{Constraint Satisfaction Problems}.
  \end{itemize}

  \begin{keyidea}
    Local search succeeds by sacrificing systematic path exploration in exchange for memory efficiency, graded optimization, and controlled exploration.
  \end{keyidea}
\end{frame}

\end{document}
